% GENERATED FILE. Edit canonical lessons, then run npm run generate:books.
% canonical-source-hash: fnv1a64:2886aa7bb87ded1a
% canonical-lessons: KA-C184-sautekayi, KA-C184-bendekayi, KA-C184-hurulikayi, KA-C184-soppu, KA-C184-elaniru, KA-R184-first-pass-words-from-chapters-176-179, KA-R184-second-pass-words-from-chapters-180-184

\chapter{Cucumber, Okra, Green beans, Leafy greens, Tender coconut water}
\label{ch:ka-cucumber-okra-green-beans-leafy-greens-tender-coconut-water}
\hlchaptermodality{184}

\begin{chapteropening}
\textbf{By the end of this chapter:} \emph{I can say a cucumber, okra, green beans, leafy greens and tender coconut water.}

Complete the last lesson of chapter 184: I can say a cucumber, okra, green beans, leafy greens and tender coconut water.
\end{chapteropening}

\section[\texorpdfstring{\kn{ಸೌತೆಕಾಯಿ} (sautekāyi)}{ಸೌತೆಕಾಯಿ (sautekāyi)}]{\kn{ಸೌತೆಕಾಯಿ} (sautekāyi) — a cucumber}

\label{lesson:KA-C184-sautekayi}



\begin{quote}
Before the new one: say the Kannada for a ladle, then the Kannada for a bucket.
\end{quote}

\subsection*{You'll want to know: \kn{ಸೌತೆಕಾಯಿ}}
\textbf{\kn{ಸೌತೆಕಾಯಿ}} — \emph{sautekāyi} — \textquotedblleft{}a cucumber\textquotedblright{}.

\begin{grammarlens}[title={What you've built}]
One more word for this chapter.
\end{grammarlens}

\subsection*{Your turn}
\begin{itemize}
\raggedright
  \item \emph{Say it:} \emph{sautekāyi}
  \item \emph{Say it:} \emph{sautekāyi}, once more
  \item \emph{Say it:} say \emph{sautekāyi}
  \item \emph{Recall:} say the Kannada for a ladle, then the Kannada for a bucket, then say \emph{sautekāyi} again
\end{itemize}

\begin{tcolorbox}[breakable,colback=teal!4,colframe=teal!35!black,title={Before you move on}]
What does \kn{ಸೌತೆಕಾಯಿ} mean? (\textquotedblleft{}A cucumber\textquotedblright{}.) Say it once more.
\end{tcolorbox}
\section[\texorpdfstring{\kn{ಬೆಂಡೆಕಾಯಿ} (beṇḍekāyi)}{ಬೆಂಡೆಕಾಯಿ (beṇḍekāyi)}]{\kn{ಬೆಂಡೆಕಾಯಿ} (beṇḍekāyi) — okra, lady's finger}

\label{lesson:KA-C184-bendekayi}



\begin{quote}
Before the new one: say the Kannada for a bucket, then the Kannada for a cucumber.
\end{quote}

\subsection*{You'll want to know: \kn{ಬೆಂಡೆಕಾಯಿ}}
\textbf{\kn{ಬೆಂಡೆಕಾಯಿ}} — \emph{beṇḍekāyi} — \textquotedblleft{}okra, lady's finger\textquotedblright{}.

\begin{grammarlens}[title={What you've built}]
One more word for this chapter.
\end{grammarlens}

\subsection*{Your turn}
\begin{itemize}
\raggedright
  \item \emph{Say it:} \emph{beṇḍekāyi}
  \item \emph{Say it:} \emph{beṇḍekāyi}, once more
  \item \emph{Say it:} say \emph{beṇḍekāyi}
  \item \emph{Recall:} say the Kannada for a bucket, then the Kannada for a cucumber, then say \emph{beṇḍekāyi} again
\end{itemize}

\begin{tcolorbox}[breakable,colback=teal!4,colframe=teal!35!black,title={Before you move on}]
What does \kn{ಬೆಂಡೆಕಾಯಿ} mean? (\textquotedblleft{}Okra, lady's finger\textquotedblright{}.) Say it once more.
\end{tcolorbox}
\section[\texorpdfstring{\kn{ಹುರುಳಿಕಾಯಿ} (huruḷikāyi)}{ಹುರುಳಿಕಾಯಿ (huruḷikāyi)}]{\kn{ಹುರುಳಿಕಾಯಿ} (huruḷikāyi) — green beans}

\label{lesson:KA-C184-hurulikayi}



\begin{quote}
Before the new one: say the Kannada for a cucumber, then the Kannada for okra.
\end{quote}

\subsection*{You'll want to know: \kn{ಹುರುಳಿಕಾಯಿ}}
\textbf{\kn{ಹುರುಳಿಕಾಯಿ}} — \emph{huruḷikāyi} — \textquotedblleft{}green beans\textquotedblright{}.

\begin{grammarlens}[title={What you've built}]
One more word for this chapter.
\end{grammarlens}

\subsection*{Your turn}
\begin{itemize}
\raggedright
  \item \emph{Say it:} \emph{huruḷikāyi}
  \item \emph{Say it:} \emph{huruḷikāyi}, once more
  \item \emph{Say it:} say \emph{huruḷikāyi}
  \item \emph{Recall:} say the Kannada for a cucumber, then the Kannada for okra, then say \emph{huruḷikāyi} again
\end{itemize}

\begin{tcolorbox}[breakable,colback=teal!4,colframe=teal!35!black,title={Before you move on}]
What does \kn{ಹುರುಳಿಕಾಯಿ} mean? (\textquotedblleft{}Green beans\textquotedblright{}.) Say it once more.
\end{tcolorbox}
\section[\texorpdfstring{\kn{ಸೊಪ್ಪು} (soppu)}{ಸೊಪ್ಪು (soppu)}]{\kn{ಸೊಪ್ಪು} (soppu) — leafy greens}

\label{lesson:KA-C184-soppu}



\begin{quote}
Before the new one: say the Kannada for okra, then the Kannada for green beans.
\end{quote}

\subsection*{You'll want to know: \kn{ಸೊಪ್ಪು}}
\textbf{\kn{ಸೊಪ್ಪು}} — \emph{soppu} — \textquotedblleft{}leafy greens\textquotedblright{}.

\begin{grammarlens}[title={What you've built}]
One more word for this chapter.
\end{grammarlens}

\subsection*{Your turn}
\begin{itemize}
\raggedright
  \item \emph{Say it:} \emph{soppu}
  \item \emph{Say it:} \emph{soppu}, once more
  \item \emph{Say it:} say \emph{soppu}
  \item \emph{Recall:} say the Kannada for okra, then the Kannada for green beans, then say \emph{soppu} again
\end{itemize}

\begin{tcolorbox}[breakable,colback=teal!4,colframe=teal!35!black,title={Before you move on}]
What does \kn{ಸೊಪ್ಪು} mean? (\textquotedblleft{}Leafy greens\textquotedblright{}.) Say it once more.
\end{tcolorbox}
\section[\texorpdfstring{\kn{ಎಳನೀರು} (eḷanīru)}{ಎಳನೀರು (eḷanīru)}]{\kn{ಎಳನೀರು} (eḷanīru) — tender coconut water}

\label{lesson:KA-C184-elaniru}



\begin{quote}
Before the new one: say the Kannada for green beans, then the Kannada for leafy greens.
\end{quote}

\subsection*{You'll want to know: \kn{ಎಳನೀರು}}
\textbf{\kn{ಎಳನೀರು}} — \emph{eḷanīru} — \textquotedblleft{}tender coconut water\textquotedblright{}.

\begin{grammarlens}[title={What you've built}]
One more word for this chapter.
\end{grammarlens}

\subsection*{Your turn}
\begin{itemize}
\raggedright
  \item \emph{Say it:} \emph{eḷanīru}
  \item \emph{Say it:} \emph{eḷanīru}, once more
  \item \emph{Say it:} say \emph{eḷanīru}
  \item \emph{Recall:} say the Kannada for green beans, then the Kannada for leafy greens, then say \emph{eḷanīru} again
\end{itemize}

\begin{tcolorbox}[breakable,colback=teal!4,colframe=teal!35!black,title={Before you move on}]
What does \kn{ಎಳನೀರು} mean? (\textquotedblleft{}Tender coconut water\textquotedblright{}.) Say it once more.
\end{tcolorbox}
\section[(dialogue)]{Words from chapters 176-179 — a first pass}

\label{lesson:KA-R184-first-pass-words-from-chapters-176-179}



\begin{quote}
Say the words for leafy greens and tender coconut water.
\end{quote}

\subsection*{Your turn}
Say these aloud:

\begin{itemize}
\raggedright
  \item for the sake of, otherwise, even so, in case, instead
  \item exactly, naturally, a result, an effect, a difference
  \item deliberately, by chance, if you want, in case, by the way
  \item fresh, stale, raw, bland, expensive
\end{itemize}

\begin{tcolorbox}[breakable,colback=teal!4,colframe=teal!35!black,title={Before you move on}]
For the sake of? (\textbf{\kn{ಸಲುವಾಗಿ}}, \emph{saluvāgi}.) Otherwise? (\textbf{\kn{ಇಲ್ಲದಿದ್ದರೆ}}, \emph{illadiddare}.) Even so? (\textbf{\kn{ಆದರೂ}}, \emph{ādarū}.) In case? (\textbf{\kn{ಒಂದು} \kn{ವೇಳೆ}}, \emph{ondu vēḷe}.) Instead? (\textbf{\kn{ಬದಲು}}, \emph{badalu}.) Exactly? (\textbf{\kn{ನಿಖರವಾಗಿ}}, \emph{nikharavāgi}.) Naturally? (\textbf{\kn{ಸಹಜವಾಗಿ}}, \emph{sahajavāgi}.) A result? (\textbf{\kn{ಫಲಿತಾಂಶ}}, \emph{phalitāṃśa}.) An effect? (\textbf{\kn{ಪರಿಣಾಮ}}, \emph{pariṇāma}.) A difference? (\textbf{\kn{ವ್ಯತ್ಯಾಸ}}, \emph{vyatyāsa}.) Deliberately? (\textbf{\kn{ಬೇಕಂತಲೇ}}, \emph{bēkantalē}.) By chance? (\textbf{\kn{ಅಕಸ್ಮಾತ್}}, \emph{akasmāt}.) If you want? (\textbf{\kn{ಬೇಕಾದರೆ}}, \emph{bēkādare}.) In case? (\textbf{\kn{ಯಾವುದಕ್ಕೂ}}, \emph{yāvudakkū}.) By the way? (\textbf{\kn{ಅಂದಹಾಗೆ}}, \emph{andahāge}.) Fresh? (\textbf{\kn{ತಾಜಾ}}, \emph{tājā}.) Stale? (\textbf{\kn{ಹಳಸಿದ}}, \emph{haḷasida}.) Raw? (\textbf{\kn{ಹಸಿ}}, \emph{hasi}.) Bland? (\textbf{\kn{ಸಪ್ಪೆ}}, \emph{sappe}.) Expensive? (\textbf{\kn{ದುಬಾರಿ}}, \emph{dubāri}.)
\end{tcolorbox}
\section[(dialogue)]{Words from chapters 180-184 — a second pass}

\label{lesson:KA-R184-second-pass-words-from-chapters-180-184}



\begin{quote}
Say the words for free of charge and government.
\end{quote}

\subsection*{Your turn}
Say these aloud:

\begin{itemize}
\raggedright
  \item free of charge, government, private, local, compulsory
  \item useful, special, favourite, modern, traditional
  \item stingy, stubborn, rotten, ripe, one's own
  \item a frying pan, a pressure cooker, a bowl, a ladle, a bucket
  \item a cucumber, okra, green beans, leafy greens, tender coconut water
\end{itemize}

\begin{tcolorbox}[breakable,colback=teal!4,colframe=teal!35!black,title={Before you move on}]
Free of charge? (\textbf{\kn{ಉಚಿತ}}, \emph{ucita}.) Government? (\textbf{\kn{ಸರ್ಕಾರಿ}}, \emph{sarkāri}.) Private? (\textbf{\kn{ಖಾಸಗಿ}}, \emph{khāsagi}.) Local? (\textbf{\kn{ಸ್ಥಳೀಯ}}, \emph{sthaḷīya}.) Compulsory? (\textbf{\kn{ಕಡ್ಡಾಯ}}, \emph{kaḍḍāya}.) Useful? (\textbf{\kn{ಉಪಯುಕ್ತ}}, \emph{upayukta}.) Special? (\textbf{\kn{ವಿಶೇಷ}}, \emph{viśēṣa}.) Favourite? (\textbf{\kn{ಅಚ್ಚುಮೆಚ್ಚಿನ}}, \emph{accumeccina}.) Modern? (\textbf{\kn{ಆಧುನಿಕ}}, \emph{ādhunika}.) Traditional? (\textbf{\kn{ಸಾಂಪ್ರದಾಯಿಕ}}, \emph{sāmpradāyika}.) Stingy? (\textbf{\kn{ಜಿಪುಣ}}, \emph{jipuṇa}.) Stubborn? (\textbf{\kn{ಹಠಮಾರಿ}}, \emph{haṭhamāri}.) Rotten? (\textbf{\kn{ಕೊಳೆತ}}, \emph{koḷeta}.) Ripe? (\textbf{\kn{ಮಾಗಿದ}}, \emph{māgida}.) One's own? (\textbf{\kn{ಸ್ವಂತ}}, \emph{svanta}.) A frying pan? (\textbf{\kn{ಬಾಣಲೆ}}, \emph{bāṇale}.) A pressure cooker? (\textbf{\kn{ಕುಕ್ಕರ್}}, \emph{kukkar}.) A bowl? (\textbf{\kn{ಬಟ್ಟಲು}}, \emph{baṭṭalu}.) A ladle? (\textbf{\kn{ಸೌಟು}}, \emph{sauṭu}.) A bucket? (\textbf{\kn{ಬಕೆಟ್}}, \emph{bakeṭ}.) A cucumber? (\textbf{\kn{ಸೌತೆಕಾಯಿ}}, \emph{sautekāyi}.) Okra? (\textbf{\kn{ಬೆಂಡೆಕಾಯಿ}}, \emph{beṇḍekāyi}.) Green beans? (\textbf{\kn{ಹುರುಳಿಕಾಯಿ}}, \emph{huruḷikāyi}.) Leafy greens? (\textbf{\kn{ಸೊಪ್ಪು}}, \emph{soppu}.) Tender coconut water? (\textbf{\kn{ಎಳನೀರು}}, \emph{eḷanīru}.)
\end{tcolorbox}
